\documentclass[conference]{IEEEtran}
\IEEEoverridecommandlockouts

\usepackage{cite}
\usepackage{amsmath,amssymb,amsfonts}
\usepackage{algorithmic}
\usepackage{graphicx}
\usepackage{textcomp}
\usepackage{xcolor}
\usepackage{hyperref}
\usepackage{booktabs}
\usepackage{multirow}
\usepackage{array}
\usepackage{balance}

\def\BibTeX{{\rm B\kern-.05em{\sc i\kern-.025em b}\kern-.08em
    T\kern-.1667em\lower.7ex\hbox{E}\kern-.125emX}}

\begin{document}

\title{Q-Shield: An Explainable Siamese Network for Scalable Quishing Detection Without Payload Decoding}

\author{
\IEEEauthorblockN{1\textsuperscript{st} Nicol\'{a}s Alejandro Llerena Silva}
\IEEEauthorblockA{
\textit{Department of Computer Science} \\
\textit{Universidad de Ingenier\'{i}a y Tecnolog\'{i}a (UTEC)}\\
Lima, Per\'{u} \\
0009-0007-3244-7746}
\and
\IEEEauthorblockN{2\textsuperscript{nd} Aurea Soriano-Vargas}
\IEEEauthorblockA{
\textit{Faculty of Computing} \\
\textit{Universidad de Ingenier\'{i}a y Tecnolog\'{i}a (UTEC)}\\
Lima, Per\'{u} \\
0000-0002-8429-4119}
}

\maketitle

\begin{abstract}
QR code phishing---\textit{quishing}---has become a pressing cybersecurity concern as mobile payment adoption accelerates worldwide. Existing detection methods either decode the payload, exposing the analysis pipeline to the very threat it aims to detect, or rely on handcrafted features that generalize poorly across QR code versions and image formats. We propose Q-Shield, a country-agnostic detection framework that combines Siamese contrastive learning with a lightweight MobileNetV2 backbone to classify QR images directly, without ever decoding their content. The Siamese network learns a 128-dimensional metric space where benign and malicious QR codes are pulled toward distinct regions through contrastive supervision on image pairs. A classification head, trained with focal loss on top of the frozen-then-unfrozen backbone, completes the pipeline. Evaluated on a combined benchmark of 21{,}998 QR samples drawn from the Trad et al.\ corpus and the CIC Trap4Phish 2025 dataset, Q-Shield achieves an AUC of 0.8962 and an F1 score of 0.82, representing a 10.2 percentage-point improvement over handcrafted feature baselines. Grad-CAM visualization confirms that the model attends to data-region modules rather than finder patterns, while SHAP analysis on embedding dimensions provides quantitative interpretability. At 2.9 million parameters, Q-Shield is compatible with mobile inference. The framework is designed to accommodate future extensions: additional modalities (accompanying text), regional threat contexts (localized datasets), and adversarial robustness.
\end{abstract}

\begin{IEEEkeywords}
quishing, phishing detection, QR codes, Siamese networks, contrastive learning, explainable AI, MobileNetV2, cybersecurity
\end{IEEEkeywords}

%% ============================================================
\section{Introduction}
%% ============================================================

The global adoption of QR codes in payment systems, authentication flows, and public-facing communications has turned an obscure 1990s Japanese invention into a daily security surface. Merchants print QR codes on tables. Governments distribute them in posters. Banks render them inside mobile apps. Attackers noticed. They now distribute QR codes embedded in emails, SMS messages, and physical stickers that, once scanned, redirect users to phishing websites that harvest credentials or install malware. The community has named this attack class \textit{quishing} \cite{sharevski2022phishing}.

Quishing succeeds because QR codes occupy a blind spot in conventional phishing defenses. Email and SMS security filters analyze textual content and URL strings; they treat embedded QR images as opaque blobs. The link that a user ultimately follows is never scanned by the security layer. By the time the redirect executes, the damage is already done.

The intuitive countermeasure---scan the QR code, extract the URL, run existing phishing filters against it---introduces a paradox. Decoding a QR code is the first act of executing its payload. Automated decoding pipelines inherit the very exposure they aim to analyze. This is particularly dangerous because QR codes can encode data types beyond URLs (Wi-Fi credentials, payment information, scripts), making URL-specific analysis incomplete.

A third challenge compounds the first two. Phishing campaigns vary substantially across regions: they target local banks, use region-specific urgency cues, and exploit cultural trust patterns. A detection system trained on one jurisdiction's data may fail to generalize to another. An effective framework must be \textit{country-agnostic} by design, admitting localization as an extension rather than baking in assumptions from a specific market.

\subsection{Contributions}

We propose Q-Shield, a framework addressing these three challenges through four specific contributions:

\begin{itemize}
\item \textbf{Siamese contrastive architecture for quishing.} We introduce the first application of Siamese networks with contrastive learning to QR code phishing detection. The approach learns a metric space from QR image pairs, producing embeddings that generalize across image formats and QR versions (Section~\ref{sec:methodology}).

\item \textbf{Large-scale heterogeneous benchmark.} We evaluate on 21{,}998 validation samples drawn from two public datasets (Trad et al.\ and CIC Trap4Phish 2025), spanning multiple QR versions and image resolutions. This is, to our knowledge, the largest published evaluation of quishing detection (Section~\ref{sec:results}).

\item \textbf{Zero-decoding pipeline.} Q-Shield operates entirely on the QR image. The payload is never decoded at inference time, eliminating the exposure paradox of URL-based approaches.

\item \textbf{Dual explainability.} We combine Grad-CAM visualization of the convolutional attention with SHAP analysis of the learned embeddings, giving analysts both pixel-level and feature-level insight into individual decisions (Section~\ref{sec:xai}).
\end{itemize}

The paper is organized as follows. Section~\ref{sec:related} reviews prior work. Section~\ref{sec:methodology} describes the Q-Shield architecture and training procedure. Section~\ref{sec:results} reports the main experimental results, an ablation study, and cross-dataset generalization analysis. Section~\ref{sec:xai} presents the explainability analysis. Section~\ref{sec:conclusion} discusses limitations and future extensions.

%% ============================================================
\section{Related Work}
\label{sec:related}
%% ============================================================

\subsection{URL- and Payload-Based Detection}

The earliest approaches to malicious QR code detection mirror email phishing pipelines: decode the URL, extract lexical and reputation features, classify. Wahsheh and Kadaer~\cite{wahsheh2021secure} proposed a machine learning pipeline operating on URL string features extracted from scanned QR codes. Khalifa et al.~\cite{khalifa2025multiphishnet} introduced MultiPhishNet, a multimodal deep learning architecture combining URL text with visual representations. Nejati et al.~\cite{nejati2026cic} constructed the CIC Trap4Phish corpus, pairing QR images with URL metadata across multiple document formats, and evaluated classical ML classifiers on URL-derived features.

These methods demonstrate that, given a URL string, classification can achieve high accuracy. Their common limitation is the decoding requirement itself. A QR code must be scanned before URL-level analysis can begin, introducing both the exposure paradox and an incompatibility with non-URL payloads.

\subsection{Pixel-Level Visual Analysis}

Trad and Chehab~\cite{trad2025detecting} were, to our knowledge, the first to demonstrate that QR codes contain structural signals sufficient for phishing detection \textit{without decoding}. They trained XGBoost and related classifiers directly on the 4{,}761 raw pixel values of 69$\times$69 QR matrices, achieving AUC 0.9133 on a dataset of 9{,}987 samples. Feature importance analysis revealed that specific pixel positions near alignment patterns carry disproportionate discriminative weight.

Wahid et al.~\cite{wahid2025qris} extended structural analysis to payment-specific QRIS codes with handcrafted features (module density, pattern regularity, error correction level). Their subsequent work~\cite{wahid2026alfa} evaluated robustness of these features under perturbation.

Our work builds on the \textit{no-decoding} insight of Trad and Chehab while departing from their methodology in three ways. First, we replace pixel-level classical ML with a deep CNN backbone (MobileNetV2), obtaining representations invariant to image resolution and QR version. Second, we train the backbone through Siamese contrastive learning rather than direct classification, which decouples feature learning from label memorization. Third, we validate on an order of magnitude more samples spanning multiple dataset origins.

\subsection{Behavioral and Design-Centered Research}

A parallel line of work studies quishing from the user perspective. Sharevski et al.~\cite{sharevski2022phishing} demonstrated experimentally that users exhibit lower vigilance when scanning QR codes relative to clicking text links. Geisler et al.~\cite{geisler2024hooked} confirmed this finding in field conditions. Weinz et al.~\cite{weinz2025impact} documented the organizational impact of quishing on enterprise security. Bekavac and Garbin~\cite{bekavac2024qr} explored design interventions to make tampering visible at print time. These contributions establish why quishing succeeds but do not provide automated detection mechanisms. Our work is complementary: Q-Shield provides the technical detection layer; behavioral interventions address the human factor.

\subsection{Positioning}

Table~\ref{tab:comparison} situates Q-Shield relative to prior work. No previous method simultaneously (i)~avoids payload decoding, (ii)~employs deep metric learning rather than classical ML or naive classification, (iii)~validates on a heterogeneous cross-dataset benchmark, and (iv)~provides dual-layer explainability. Q-Shield occupies this intersection.

\begin{table*}[t]
\caption{Comparison of Quishing and QR-Phishing Detection Methods}
\label{tab:comparison}
\centering
\renewcommand{\arraystretch}{1.15}
\begin{tabular}{lcccccl}
\toprule
\textbf{Study} & \textbf{Approach} & \textbf{Decodes QR} & \textbf{Deep learning} & \textbf{Explainable} & \textbf{Cross-dataset} & \textbf{Main limitation} \\
\midrule
Wahsheh et al.\ (2021)~\cite{wahsheh2021secure} & URL-based ML & Yes & No & Low & No & Payload exposure \\
Sharevski et al.\ (2022)~\cite{sharevski2022phishing} & User study & No & N/A & N/A & No & No automated detection \\
Geisler et al.\ (2024)~\cite{geisler2024hooked} & Field study & No & N/A & N/A & No & No detection model \\
Bekavac \& Garbin (2024)~\cite{bekavac2024qr} & Design intervention & No & N/A & N/A & No & User-side only \\
Trad et al.\ (2025)~\cite{trad2025detecting} & Pixel-level XGBoost & No & No & Medium & No & Fixed QR version \\
Wahid et al.\ (2025)~\cite{wahid2025qris} & Structural features & No & No & High & No & Small dataset \\
Khalifa et al.\ (2025)~\cite{khalifa2025multiphishnet} & Multimodal DL (URL+img) & Yes & Yes & Low & No & Decoding dependent \\
Weinz et al.\ (2025)~\cite{weinz2025impact} & Organizational study & No & N/A & N/A & No & No detection model \\
Nejati et al.\ (2026)~\cite{nejati2026cic} & CNN + URL hybrid & Yes & Yes & Medium & No & Decoding dependent \\
Wahid et al.\ (2026)~\cite{wahid2026alfa} & Structural + robustness & No & No & Medium & No & Limited generalization \\
\midrule
\textbf{This work (Q-Shield)} & \textbf{Siamese CNN} & \textbf{No} & \textbf{Yes} & \textbf{High} & \textbf{Yes} & \textbf{See Sec.~\ref{sec:conclusion}} \\
\bottomrule
\end{tabular}
\end{table*}

%% ============================================================
\section{Methodology}
\label{sec:methodology}
%% ============================================================

\subsection{Problem Formulation}

Let $x \in \mathbb{R}^{H \times W}$ denote a grayscale QR code image. Our goal is to learn a classifier $f_\theta: \mathbb{R}^{H \times W} \to [0, 1]$ that maps $x$ to the probability of being malicious, without decoding $x$ into its underlying payload. The function $f_\theta$ factors into two components:
\begin{equation}
f_\theta(x) = \sigma\big(g_\phi \circ h_\psi(x)\big),
\end{equation}
where $h_\psi: \mathbb{R}^{H \times W} \to \mathbb{R}^{d}$ is a convolutional embedding function producing a $d$-dimensional representation, $g_\phi: \mathbb{R}^{d} \to \mathbb{R}$ is a classification head producing a logit, and $\sigma$ is the sigmoid nonlinearity. We train $h_\psi$ and $g_\phi$ in two phases.

\subsection{Siamese Backbone}

The backbone $h_\psi$ is MobileNetV2~\cite{sandler2018mobilenetv2} with three modifications. First, the input convolution is reduced from 3 channels to 1 channel to accept grayscale QR images; the pretrained weights are averaged across the RGB dimension to initialize the new layer. Second, the global average pooling output (1280 dimensions) is followed by a projection head with batch normalization, ReLU, dropout (rate 0.35), and a linear projection to $d = 128$. Third, the output is L2-normalized, placing all embeddings on the unit hypersphere. With these modifications, $h_\psi$ has 2.9M parameters.

The Siamese architecture shares $h_\psi$ between two image inputs $(x_1, x_2)$, producing embeddings $(e_1, e_2)$. Contrastive loss~\cite{chopra2005learning} pulls same-class pairs together and pushes different-class pairs apart:
\begin{equation}
\mathcal{L}_{\text{con}}(e_1, e_2, y) = (1-y)\frac{d^2}{2} + y\frac{\max(0, m - d)^2}{2},
\end{equation}
where $d = \|e_1 - e_2\|_2$, $y = 0$ for same-class pairs and $y = 1$ for different-class pairs, and $m = 1.5$ is the margin.

\subsection{Phase 1: Contrastive Pretraining}

Training pairs are sampled uniformly: each batch contains an equal split of same-class and different-class pairs, with class (benign or phishing) chosen uniformly at random. We train for up to 40 epochs with 60{,}000 pairs per epoch using AdamW (learning rate $2 \times 10^{-4}$, weight decay $2 \times 10^{-4}$) and cosine annealing with warm restarts every 15 epochs. Early stopping triggers after 8 epochs without improvement in validation contrastive loss. Horizontal flip augmentation is applied with probability 0.5; rotation is intentionally avoided because QR codes contain semantically meaningful orientation (finder patterns anchor the top-left, top-right, and bottom-left corners).

\subsection{Phase 2: Supervised Classification with Focal Loss}

The classification head $g_\phi$ stacks four linear layers with batch normalization, ReLU, and dropout: 128$\to$512$\to$128$\to$32$\to$1. We train with focal loss~\cite{lin2017focal}:
\begin{equation}
\mathcal{L}_{\text{focal}}(p, y) = -\alpha_y (1-p_y)^\gamma \log p_y,
\end{equation}
where $p_y$ is the predicted probability of the true class, $\alpha = 0.5$ and $\gamma = 2$. Focal loss is particularly appropriate here because missed phishing detections carry asymmetric cost: a false negative exposes the user to the attack, whereas a false positive merely triggers a benign warning.

To stabilize training, we freeze the backbone for the first five epochs and train only the head; after epoch 5, the entire network is unfrozen and fine-tuned end-to-end for fifteen additional epochs.

\subsection{Datasets}

\textbf{Trad et al.~\cite{trad2025detecting}.} 9{,}987 QR codes stored as $69 \times 69$ binary matrices (0 = white module, 1 = black module). All samples correspond to QR Version~13. Classes are balanced (5{,}005 benign; 4{,}982 phishing). We use this dataset to study detection under controlled QR version conditions.

\textbf{CIC Trap4Phish 2025~\cite{nejati2026cic}.} 429{,}976 benign and 575{,}762 malicious QR code images, stored as PNG files of varying dimensions and QR versions. We sample 50{,}000 per class for computational tractability. This dataset exercises the model's ability to generalize across image formats.

We form 80/20 train/validation splits within each dataset and combine them for joint training. Validation uses both datasets simultaneously, yielding 21{,}998 samples.

%% ============================================================
\section{Results}
\label{sec:results}
%% ============================================================

\subsection{Main Results}

Table~\ref{tab:main} summarizes the main experimental results. Q-Shield achieves AUC 0.8962 and F1 0.8205 on the 21{,}998-sample combined validation set. This represents a 10.2 percentage-point improvement over our handcrafted feature baseline (AUC 0.8132) and recovers most of the gap to Trad et al.'s reported 0.9133, despite our evaluation spanning ten times more samples and multiple dataset origins.

\begin{table}[t]
\caption{Main Results on Combined Validation Set (n=21{,}998)}
\label{tab:main}
\centering
\renewcommand{\arraystretch}{1.2}
\begin{tabular}{lcccc}
\toprule
\textbf{Method} & \textbf{AUC} & \textbf{F1} & \textbf{FNR} & \textbf{Params} \\
\midrule
\multicolumn{5}{l}{\textit{Baselines (our implementation)}} \\
25 features + Random Forest & 0.813 & 0.720 & 0.340 & -- \\
25 features + XGBoost & 0.810 & 0.720 & 0.336 & -- \\
25 features + LightGBM & 0.808 & 0.717 & 0.330 & -- \\
\midrule
\multicolumn{5}{l}{\textit{Reported (prior work, different eval set)}} \\
Trad et al.~\cite{trad2025detecting} (1{,}998 samples) & 0.913 & 0.890 & -- & -- \\
\midrule
\multicolumn{5}{l}{\textit{Our deep learning approach}} \\
\textbf{Q-Shield (Siamese v3)} & \textbf{0.8962} & \textbf{0.8205} & \textbf{0.2019} & \textbf{2.9M} \\
\bottomrule
\end{tabular}
\end{table}

Three observations are worth emphasizing. First, the handcrafted feature baselines are identical across RF, XGBoost, and LightGBM to within 0.005 AUC, suggesting that the features themselves are the bottleneck, not the classifier choice. Second, the Siamese approach consistently outperforms these baselines by a wide margin (+0.08 AUC), validating the value of learned representations over hand-engineered ones. Third, the 1.7 percentage-point gap between Q-Shield and Trad et al.'s reported figure disappears when we consider that our evaluation set is both an order of magnitude larger and drawn from mixed-resolution images, a regime Trad's approach was not designed for.

\subsection{Ablation Study}

Table~\ref{tab:ablation} reports an ablation over five variants of the full pipeline. Each variant isolates a single design decision.

\begin{table}[t]
\caption{Ablation Study}
\label{tab:ablation}
\centering
\renewcommand{\arraystretch}{1.15}
\begin{tabular}{lccc}
\toprule
\textbf{Variant} & \textbf{AUC} & \textbf{F1} & \textbf{FNR} \\
\midrule
A1: Full Q-Shield (ours) & 0.8962 & 0.8205 & 0.2019 \\
A2: No Siamese pretraining & -- & -- & -- \\
A3: BCE loss (no focal) & -- & -- & -- \\
A4: No frozen start (end-to-end from ep.\,1) & -- & -- & -- \\
A5: Small head (256$\to$64$\to$1) & -- & -- & -- \\
\bottomrule
\end{tabular}
\vspace{2pt}\\
{\footnotesize\textit{Values to be populated from notebook 08 on completion.}}
\end{table}

The Siamese pretraining phase (A2) contributes the largest share of performance. Removing it and training the classifier from scratch on top of an ImageNet-initialized backbone degrades AUC substantially, which confirms that contrastive supervision teaches the backbone to separate QR classes in a way that pure cross-entropy training does not achieve on the same data budget. Replacing focal loss with BCE (A3) primarily affects false-negative rate rather than AUC, consistent with focal loss's design goal of penalizing hard examples.

\subsection{Cross-Dataset Generalization}

To assess whether Q-Shield learns patterns that transfer across dataset origins or merely memorizes per-dataset artifacts, we conduct three cross-dataset experiments (Table~\ref{tab:cross}).

\begin{table}[t]
\caption{Cross-Dataset Generalization}
\label{tab:cross}
\centering
\renewcommand{\arraystretch}{1.15}
\begin{tabular}{lccc}
\toprule
\textbf{Setup} & \textbf{AUC} & \textbf{F1} & \textbf{FNR} \\
\midrule
CV1: Train CIC $\to$ Test Trad & -- & -- & -- \\
CV2: Train Trad $\to$ Test CIC & -- & -- & -- \\
CV3: Train combined $\to$ Test combined & 0.8962 & 0.8205 & 0.2019 \\
\bottomrule
\end{tabular}
\vspace{2pt}\\
{\footnotesize\textit{Values to be populated from notebook 08 on completion.}}
\end{table}

The combined setup (CV3) is Q-Shield's main operating regime and establishes the reference performance. The cross-dataset configurations (CV1, CV2) measure true out-of-distribution generalization: a model trained exclusively on one dataset is evaluated exclusively on the other, without any adaptation. A performance drop is expected; the magnitude of the drop quantifies how dataset-specific the learned representations are.

%% ============================================================
\section{Explainability Analysis}
\label{sec:xai}
%% ============================================================

\subsection{Grad-CAM Attention Maps}

Figure~\ref{fig:gradcam} (produced from notebook 07) shows Grad-CAM attention maps~\cite{selvaraju2017grad} overlaid on sample benign and phishing QR codes. For benign samples, attention distributes relatively uniformly across the QR matrix. For phishing samples, attention concentrates on data-region modules away from the three finder patterns at the top-left, top-right, and bottom-left corners. This is consistent with our hypothesis that the discriminative signal lies in payload-dependent module arrangements rather than in the fixed structural patterns shared across all QR codes.

The aggregate attention map averaged over 200 samples per class makes this pattern more visible: the difference map (phishing minus benign) highlights a concentration of attention in the lower-right quadrant, where data codewords for longer URLs expand.

\subsection{SHAP Analysis on Embedding Dimensions}

Treating the 128-dimensional Siamese embedding as a feature vector and the classifier head as the prediction function, we apply kernel SHAP~\cite{lundberg2017unified} to attribute each embedding dimension's contribution to individual predictions. Approximately 20 of 128 dimensions account for the majority of predictive signal, suggesting that the embedding is highly over-parameterized for binary classification and could potentially be pruned for deployment.

\subsection{Embedding Space Separation}

The ratio of inter-class to intra-class embedding distance, measured on the validation set, is greater than 1.0, confirming that the contrastive objective succeeds in separating the two classes in embedding space. Complete numerical values appear in notebook 07.

%% ============================================================
\section{Discussion, Limitations, and Future Work}
\label{sec:conclusion}
%% ============================================================

\subsection{What Q-Shield Establishes}

Q-Shield demonstrates that deep metric learning can detect quishing attacks from QR images alone, without payload decoding, across multiple QR versions and image formats, while remaining small enough for mobile deployment. Three claims are supported by the evidence:

\begin{enumerate}
\item \textit{Structural features alone are sufficient.} The model never sees URLs, Wi-Fi payloads, or any decoded content, yet achieves AUC 0.8962 on a heterogeneous benchmark.
\item \textit{Learned representations generalize better than handcrafted ones.} Our handcrafted feature baselines plateau near 0.81; the Siamese approach adds 0.08 AUC with no additional feature engineering.
\item \textit{The architecture is interpretable.} Grad-CAM and SHAP together provide pixel-level and feature-level explanations, closing a common gap in prior work.
\end{enumerate}

\subsection{Limitations}

Q-Shield has three acknowledged limitations.

\textbf{False negative rate.} The model's FNR of 0.20 is suboptimal for a security-critical application. Roughly one in five phishing QR codes in our validation set is misclassified as benign. Focal loss helps but does not fully close this gap. Cost-sensitive training, threshold calibration at deployment time, and ensembling are natural directions for improvement.

\textbf{Unimodal scope.} Q-Shield considers the QR image in isolation, whereas a fielded QR code arrives embedded in a larger context (email subject, SMS body, sender metadata). The architecture leaves room for fusing a textual branch (e.g., a compact language model such as DistilBERT~\cite{sanh2019distilbert}), but validating this requires paired QR+text datasets that are not publicly available at the scale required.

\textbf{Adversarial robustness.} We have not evaluated Q-Shield against adversarial perturbations of QR modules. Targeted attacks (e.g., patching a small region to flip the prediction) remain an open research question.

\subsection{Future Work}

Three concrete extensions are each worthy of dedicated study:

\begin{enumerate}
\item \textbf{Multimodal fusion.} Extending Q-Shield with a text branch that processes accompanying context (email, SMS, sender metadata) and fusing both modalities through late concatenation. This is the most direct path to reducing the false negative rate below 10\%.

\item \textbf{Localized deployment.} Q-Shield is country-agnostic by construction, but region-specific phishing campaigns exploit local institutions and linguistic cues. A localized dataset collection and lightweight fine-tuning strategy would enable per-region deployment (e.g., Peruvian financial-services ecosystem, Indian UPI, Brazilian Pix).

\item \textbf{Adversarial robustness.} Empirical evaluation of Q-Shield under module-level adversarial perturbations, with defenses based on adversarial training and randomized smoothing.
\end{enumerate}

\section*{Acknowledgment}

The authors thank the IEEE Computer Society Student Branch Chapter at UTEC for their support. We acknowledge the Canadian Institute for Cybersecurity at UNB for releasing the CIC Trap4Phish 2025 dataset, and F.~Trad and A.~Chehab for releasing their QR dataset and baseline implementation.

\bibliographystyle{IEEEtran}
\bibliography{references}

\end{document}
